\documentclass{article}
\usepackage{graphicx}
\usepackage{caption}
\usepackage{subcaption}
\usepackage{geometry}
\geometry{margin=1in}

\title{Comparison of Tabular Q-Learning and Deep Q-Networks (DQN) for Snake}
\author{Your Name}
\date{\today}

\begin{document}

\maketitle

\section{Introduction}
Reinforcement learning (RL) has been successfully applied to classic games such as Snake. In this project, we compare two RL approaches: Tabular Q-learning and Deep Q-Networks (DQN). We analyze their performance and learning characteristics on the Snake game environment.

\section{Tabular Q-Learning}
Tabular Q-learning is a model-free RL algorithm that maintains a table of Q-values for each state-action pair. It is suitable for environments with small, discrete state spaces. In our implementation, the state was discretized, and the agent learned optimal actions through repeated interactions.

\begin{figure}[h!]
    \centering
    \includegraphics[width=0.7\textwidth]{image1.png}
    \caption{Performance of Tabular Q-learning on Snake.}
    \label{fig:tabular}
\end{figure}

\section{Deep Q-Network (DQN)}
DQN extends Q-learning by using a neural network to approximate the Q-value function, enabling it to handle high-dimensional or continuous state spaces. In our project, DQN was trained using a feature-based state representation.

\begin{figure}[h!]
    \centering
    \includegraphics[width=0.7\textwidth]{image2.png}
    \caption{Performance of DQN on Snake.}
    \label{fig:dqn}
\end{figure}

\section{Comparison}
To visually compare the learning curves of Tabular Q-learning and DQN, we present both results side by side in Figure~\ref{fig:comparison}. After training both agents for 130 episodes, we observed significant performance differences:

\begin{itemize}
    \item \textbf{Tabular Q-Learning}: Achieved a peak score of \textbf{18}
    \item \textbf{DQN}: Achieved a peak score of \textbf{54}
\end{itemize}

DQN significantly outperformed Tabular Q-learning, achieving nearly 3x higher peak scores. This demonstrates the advantage of using neural networks to approximate Q-values, allowing the agent to generalize better from the feature-based state representation compared to the discrete state space limitations of tabular methods.

\begin{figure}[h!]
    \centering
    \begin{subfigure}[b]{0.45\textwidth}
        \centering
        \includegraphics[width=\textwidth]{image1.png}
        \caption{Tabular Q-learning}
    \end{subfigure}
    \hfill
    \begin{subfigure}[b]{0.45\textwidth}
        \centering
        \includegraphics[width=\textwidth]{image2.png}
        \caption{DQN}
    \end{subfigure}
    \caption{Comparison of Tabular Q-learning and DQN results on Snake.}
    \label{fig:comparison}
\end{figure}

\section{Conclusion}
This report compared Tabular Q-learning and DQN for the Snake game. DQN outperformed Tabular Q-learning in terms of learning speed and final performance. Further work could include evaluating Convolutional DQN (ConvDQN) with raw pixel input.

% Placeholder for ConvDQN results
% \section{Convolutional DQN (ConvDQN)}
% Results and analysis to be added in future work.

\end{document} 